\partopening{Options and Constituent Information}{part:options}
{What additional contracts reveal, and which valuation assumptions make that possible.}

\section{Option prices do not have a unique bubble interpretation}\label{sec:optionrules}
The option results explain how additional contracts can constrain a
valuation amount. The practical difficulty appears in the observed quotes:
on the sampled end-of-day SPY chains, the largest-strike calls at one to
three months sit at the one-cent floor on most dates, at median strikes
1.07--1.41 times spot. The available far wing is therefore measured only to
the tick. Reading such quotes as a bubble bound can produce a very small
number largely determined by the pricing convention and quote resolution.
\source{experiment 014}

The central issue is whether the observed call represents an expectation
price, a parity-restoring price, or an exercise value. Assuming that a call
carries the defect and then reading a tiny quote as evidence of its absence
makes the conclusion depend on the premise being assessed. The results below
are consequently valuation identities and bounds under named conventions;
the research did not obtain a usable market bubble test from these quotes.
The algebra is nevertheless informative about which observations would
constrain a defect and which can leave it hidden.

Normalize spot to one, take zero rates and dividends, and fix a common
pricing law with $m=\E_Q[X_T]$ and $\delta=1-m$. European expectation
prices satisfy
\begin{equation}\label{eq:parityexpect}
 c(K)=\E_Q[(X_T-K)^+],\quad p(K)=\E_Q[(K-X_T)^+],
 \qquad c(K)-p(K)=m-K.
\end{equation}
Thus $1-K-c+p=\delta$. Under another declared rule, leave the put at
$p$ and set the call to $C=c+\delta$. Then $1-K-C+p=0$ even when
$\delta>0$. This second rule restores ordinary parity.

The distinction builds on the strict-local-martingale option literature.
Cox and Hobson explain how parity and exercise values depend on the
admissible trading strategies \cite{cox2005local}.
Carr, Fisher, and Ruf obtain parity restoration from their
joint-superreplication setting \cite{carr2012hedging}; it is not a
universal conclusion that market calls follow that rule. The lab uses the
two rules as explicit alternatives and works out their implications for
inference. Calling both simply ``risk-neutral option prices'' would hide
the assumption that determines the answer.

\subsection{Exercise value and the parity residual}
For the corresponding continuously exercisable American call, the lab's
localization argument gives $A_Q\geq c+\delta$. Write
$A_Q=c+\delta+u_A$, where $u_A\geq0$ denotes any excess over that lower
bound. With a European expectation put in every row:
\begin{table}[htbp]\centering\small
\begin{tabularx}{\linewidth}{@{}YllY@{}}\toprule
Call convention & Call value & $1-K-C+p$ & Reading\\
\midrule
European expectation & $c$ & $\delta$ & Residual equals the defect\\
European parity-restoring & $c+\delta$ & $0$ & A defect can be hidden by parity\\
American exercise & $c+\delta+u_A$ & $-u_A$ & Residual concerns exercise excess\\
\bottomrule\end{tabularx}
\caption{Residuals under declared valuation rules. The table is not a
classification algorithm for market quotes.}\label{tab:operators}
\end{table}

Contract metadata distinguish American from European exercise even if their
prices coincide. Prices alone need not distinguish valuation mechanisms:
on the equality face $u_A=0$, the parity-restored European value and the
American value agree. Replacing the put by an American put adds its own
exercise premium and changes the table. Neither a zero residual nor an
American premium alone identifies a defect without the rest of the model.

With simultaneous intervals for the relevant operator prices, one can
propagate errors through the residual and test the zero-residual
restriction. If the total call-plus-put error is bounded by $h$, a
positive residual exceeding $2h$ in the expectation model is sufficient
for rejection by a conservative lower-bound rule. That is a conditional
separation statement. A bid--ask spread by itself does not establish the
required coverage of a specified pricing operator.
\source{claims 014 and 017; the D31 option-observables design. This
design contains no real-data option test}

\section{Which option observations are informative?}\label{sec:optioninfo}
\subsection{The lowest strike matters}
Suppose only expectation-priced calls are observed at finitely many
strikes $K_1<\cdots<K_r$, all at or above normalized spot one. If their
prices are compatible with a nonnegative terminal law of mean at most
one, then they are also compatible with a terminal law of mean exactly
one. They cannot exclude zero defect.

The construction is simple. Starting from a compatible random variable
$X$ with mean $m<1$, compare it with $Y=\max(X,K_1)$. Every observed
call payoff is unchanged because the added mass lies below its strike.
But $\E[Y]=K_1+c(K_1)\geq1$. Mixing the two terminal laws raises the
mean to one while preserving the entire observed call panel.

This is a statement about terminal laws, not a construction of a diffusion
for every feasible law. It still makes a useful observation-design point:
adding ever-higher call strikes does not automatically identify the mean
missing below the lowest strike. An in-the-money call can exclude mean
one in this model. A matched underlying call and put identify the mean
directly through \eqref{eq:parityexpect}, conditional on expectation pricing.

The lab also constructs computable outer intervals for the defect using
call monotonicity, convex secant slopes, and quote intervals. Those
inequalities alone are not always a sharp characterization of attainable
terminal laws. Inconsistent observations must be reported as model
incompatibility, not automatically as positive bubble evidence.
\source{claim 021}

\subsection{Do options on leveraged funds add information?}
For a normalized fund, the same expectation-pricing identity yields its
defect from a call--put pair. Finite-grid inheritance then links a
positive fund defect to the underlying defect. For exposures
$0<\pi\leq1$, the quantitative link uses the factor $\pi^N$; for
liquidated long leverage it uses $\pi^N$ in the opposite inequality.
Both become weak for monetary estimation at realistic numbers of resets.

For example, $0.75^{252}$ is about $3\times10^{-32}$. An upper bound
on the fund defect divided by this factor is normally uninformative
about the underlying's defect. On the leveraged side, a positive fund
defect divided by $\pi^N$ supplies a negligible lower bound. Zero
exclusion can remain mathematically meaningful even when amount
estimation is weak, but its practical power then depends on resolving a
small residual amid costs, tracking, and quote uncertainty.

If an exact expectation-priced underlying call--put pair is already
available, it identifies the defect without fund options. Adding fund
options cannot improve on exact identification. With only underlying
out-of-the-money calls, a positive fund residual can add a restriction,
but no empirical gain specific to leverage has been demonstrated. The
inverse liquidation model in Part~\ref{part:funds} has zero fund defect
even over a strict underlying, so it cannot supply the same zero-exclusion
logic.
\source{claims 014, 018, and 021; no real-data fund-option detection
result is claimed}

\subsection{Continuous exercise creates another limiting distinction}
Let $G_t=\bar S_t+D_t$ be a finite continuous nonnegative
$Q$-local martingale, where $\bar S_t=d_tS_t\geq0$ and $d_t$ is a
positive deterministic continuous discount factor, bounded on $[0,T]$.
The cumulative discounted distributions $D$ are adapted, nonnegative,
nondecreasing and c\`adl\`ag, with $D_0=0$. A distribution jump may be
offset by a price drop while the gain process remains continuous. Write
$b_G(T)=G_0-\E_Q[G_T]$, and define
\[
 C_E(K,T)=\E_Q[(\bar S_T-d_TK)^+],\qquad
 A_Q(K,T)=\sup_{\tau\leq T}\E_Q[(\bar S_\tau-d_\tau K)^+],
\]
where the supremum includes every stopping time in $[0,T]$.
For this continuously exercisable American call,
localization at high gain levels and bounded convergence of the capped
payoff component give
\begin{equation}\label{eq:american}
 A_Q(K,T)\geq b_G(T)+C_E(K,T)\geq b_G(T).
\end{equation}
The relation between American exercise and bubbles has established
antecedents \cite{bayraktar2011strict}; it is not presented as a new
general pricing principle. The lab's treatment includes distributions
and the transfer to fixed-basket bounds.

The contrast with finite exercise is especially revealing. European
expectation calls tend to zero as $K\to\infty$. The same is true when
exercise is restricted to any fixed finite set of dates, because the
exercise payoff is bounded by a finite sum of integrable call payoffs.
Such prices therefore cannot uniformly bound a positive defect from
above as the strike increases. Continuous exercise permits stopping at
rare high levels in a way that a fixed finite exercise schedule does not.
The order of these limiting operations matters.

A tiny quoted call is consequently not an assumption-free certificate
of a tiny bubble. To use \eqref{eq:american}, the quote must provide a
justified upper bound on the specified continuous-exercise value. The
model of exercise and its relation to observed market prices are
substantive premises.

\section{From an index to its constituents}\label{sec:baskets}
For fixed positive holdings $n_i$ and a common pricing measure, let
$I_t=\sum_i n_iS_{i,t}$, $\omega_i=n_iS_{i,0}/I_0$, and
$b_i(T)=S_{i,0}-\E_Q[S_{i,T}]$. Linearity gives
\begin{equation}\label{eq:basketadd}
 b_I(T)=\sum_i n_i b_i(T),\qquad 0\leq n_i b_i(T)\leq b_I(T).
\end{equation}
This is a consequence of the definition, not a new diversification
principle. It concerns fixed holdings under one measure; a rebalanced
index needs its trading rule specified separately.

\subsection{Ordinary weights can hide all of the effect}
Suppose two names are proportional strict prices, with equal initial
weights. Their ordinary cap weights remain exactly one-half. Both can
have positive defects while neither weight declines. Thus a proposed
attribution based solely on expected losses of ordinary weights cannot
work in general.

The corrected representation uses the index share probability. Its
defining identity carries an explicit survival event up to a lifetime
$\tau^I$ \cite{kardaras2015strict}. Let
$\mu_i(t)=n_iS_{i,t}/I_t$ and
$\mu_i^\dagger(t)=\mu_i(t)\ind_{\{t<\tau^I\}}$. Then
\begin{equation}\label{eq:killed}
 \omega_i-\E_{\widetilde Q_I}[\mu_i^\dagger(T)]
       =\frac{n_i b_i(T)}{I_0},\qquad
 \sum_i\mu_i^\dagger(T)=\ind_{\{T<\tau^I\}}.
\end{equation}
These survival-adjusted, or killed, weights need not sum to one after
the lifetime. Ordinary weights do, which is why they cannot all be
supermartingales with a positive total expected loss under a single
probability. In the proportional example, all the loss is carried by
the survival factor.

The associated index hazard is also a survival-weighted mixture. If
$p_i(t)=\E_Q[S_{i,t}]/S_{i,0}$ are absolutely continuous with
$h_i=-p_i'/p_i$, then
\begin{equation}\label{eq:hazard}
 p_I=\sum_i\omega_i p_i,
 \qquad h_I=\sum_i\frac{\omega_i p_i}{p_I}h_i.
\end{equation}
The weights are not simply today's or the initial cap weights. These
are auxiliary-measure lifetimes and hazards, not directly observed
physical crash probabilities.
\source{claim 010}

\subsection{A conditional recovery formula}
Let $\nu_T$ be the pricing distribution of the terminal index. The recovery result takes this distribution as a supplied population
input. Let
$m_i(x)=\E_P[\mu_i(T)\mid I_T=x]$. If the terminal change of measure
has the form $\mathrm{d}Q/\mathrm{d}P=k_T(I_T)$, then conditioning on
the terminal index preserves the conditional mean weights. Hence
\begin{equation}\label{eq:recovery}
 \frac{n_i b_i(T)}{I_0}
   =\omega_i-\frac1{I_0}\int x m_i(x)\nu_T(\mathrm{d}x).
\end{equation}
The weaker condition that the conditional mean weights agree under
$P$ and $Q$ already suffices. The density condition is one way to imply
it, and is much stronger than loosely saying that the pricing kernel
``depends on the market.''

The formula is an identification result under that restriction, not an
estimator from one historical stock panel and today's finite option
chain. The required conditional means and pricing distribution are
population objects. The lab's estimator assumes independent physical
terminal panels and an independent supplied pricing sample; it does
not construct those inputs from overlapping historical windows.

There is a necessary compatibility check:
$0\leq\int x m_i(x)\nu_T(\mathrm{d}x)\leq n_iS_{i,0}$ for every
name. The estimated contributions summing to the index contribution
is not a check, because the identity follows automatically when the
conditional weights sum to one. This distinction is useful whenever
an attribution method appears to validate itself by matching a total.
\source{claims 016 and 044}

\subsection{What the synthetic estimation says}
The recovery estimator preserves the total on every sample and is
consistent under the stated sampling and approximation assumptions.
A histogram-based necessary-compatibility test has a conditional
finite-sample error bound with a supplied smoothness constant. In
controlled terminal-law examples, finite-sample tests detect large
violations; the fixed partitions used there do not implement an asymptotic
power construction.

In the more process-like basket example, however, the bubble contribution
is about $0.0041$ while the fixed-partition estimator bias is about
$0.0035$. Increasing the sample size does not remove fixed-range and
fixed-partition bias. No compatibility rejection was observed in the
basket cases tested. A numerical smoothness estimate gives convenient
intervals but lacks a rigorous bound; the available rigorous bound is too
large to make those intervals useful. The population identity is stronger
than this estimator's demonstrated performance.
\source{experiment 046}

\subsection{What a basket budget can and cannot identify}
Suppose an option on a tracking fund supplies a justified upper bound
$U$ on its American exercise value, and $\eta$ bounds the expected
tracking discrepancy under $Q$ relative to a frozen basket. With the
appropriate scaling $a$, the total defect is bounded by
$B_{\mathrm{cap}}=aU+\eta$. Then
\begin{equation}\label{eq:budget}
 \frac{b_i(T)}{S_{i,0}}
    \leq\frac{B_{\mathrm{cap}}}{\omega_i I_0}.
\end{equation}
The same absolute budget can give an informative relative bound for
a large constituent and a vacuous one for a tiny constituent. Tracking
under the physical measure does not supply $\eta$, and fees and basket
reconstitution can matter more than a small quoted wing call.

An upper budget alone never forces any particular name to have a
positive defect. Combined with individual intervals and a positive
lower bound on the total, the lab derives exact coordinate bounds for
the resulting box intersected by a weighted-sum constraint. That
finite-dimensional calculation is useful bookkeeping, but it cannot
create the missing monetary lower bound from a binary price-based
signal. A small positive upper bound is a tolerance certificate, not
proof of exact absence.
\source{claims 017 and 019}

\section{What the four parts establish}\label{sec:remaining}
The research gives a connected account of three practical objects: a
price--variance signal, the wealth produced by a trading rule, and an
expectation shortfall under a pricing law. The daily statistic detects
specified level alternatives substantially more often than calibrated
GSADF and identifies a coherent early pattern across the four homebuilders.
The intraday score adds drawdown information to the explosive-root
benchmark. Calibration determines the reference against which either
pattern is unusual.

The theory explains why those measurements and a valuation defect are
related but distinct. Rapidly growing local volatility and the positive
shock alignment in Sin's model provide two mechanisms. The exposure
classification shows how rebalancing can change strictness, and the
finite-grid and liquidation results explain why actual contract rules
matter. A bounded liquidated inverse fund and an unbounded leveraged-long
fund can have different expectation properties over the same underlying.

Options and constituent data then put structure on the valuation question.
A matched expectation-priced call and put determine a terminal mean;
a high-strike call panel alone can leave it unidentified. A fixed basket
adds constituent defects, while attribution needs more information than
ordinary weights. These results give the practitioner useful distinctions:
which pattern a monitor measures, which exposure a contract creates, and
which assumptions make a monetary bound informative.
